% !PRINTED 55
\weberpage{55}
%<<CONT
same, and may be called the divisor of the class. If $d, d'$ are the divisors of two classes $a, a'$, then the greatest common divisor of $m$ and $dd'$ is the divisor of the class $aa'$. From this it follows that the number classes whose numbers are relatively prime to $m$ always yield such number classes again under composition, and for these the requirement of \S.\,1, 3. is then also fulfilled.

1.\ \gesperrt{The number classes whose individuals are relatively prime to the modulus thus form, under composition by multiplication, an Abelian group.}

This group is the object of our considerations, in which our chief goal shall be the determination of a basis.

We shall denote by $N$ every number class that contains only numbers coprime to $m$, and the group of the number classes $N$, whose existence we have now demonstrated, by $\mathfrak{N}$. By $n$ we shall denote every number coprime to $m$.

The degree of this group is as great as the number of the relative primes to $m$ that are at the same time positive and not greater than $m$, and this number we have determined in \S.\,132 of the first volume. It is

\weq{2}{\mu = \varphi(m) = 2^{\lambda-1} q_1^{\varkappa_1 - 1}(q_1 - 1)\, q_2^{\varkappa_2 - 1}(q_2 - 1) \ldots}

or

\[
= q_1^{\varkappa_1 - 1}(q_1 - 1)\, q_2^{\varkappa_2 - 1}(q_2 - 1) \ldots,
\]

when $\lambda = 0$.

The degree of an element $N$ of this group is the least positive exponent $e$ to which one must raise a representative $n$ of $N$ in order that $n^e$ become congruent to unity with respect to the modulus $m$. Since every $e$ must be a divisor of the degree of the group $\varphi(m)$, there follows the \gesperrt{generalised Fermat's theorem}:

\weq{3}{n^{\varphi(m)} \equiv 1\ (\mathrm{mod}\ m),}

a formula which holds for every number $n$ coprime to $m$.

If now $q$ is one of the odd prime factors of $m$, and $q^\varkappa$ the highest power of $q$ dividing $m$, then we take a primitive root $g$ of $q$, which, if $\varkappa$ is greater than 1, we choose so that $g^q - g$ is not divisible by $q^2$ (Bd. I, \S.\,184). For the sake of brevity we shall call a number $g$ satisfying this last condition a \gesperrt{primitive root of $q^2$}.
